\documentclass[11pt]{scrartcl}

\usepackage[stixtwo]{fontsetup}
\usepackage{amsmath}

\begin{document}
A rendőrelv (bármelyik változata elfogadható, akár a speciális is): ha
$a_n \le b_n \le c_n$ elég nagy $n$-re, és $a_n \to A$, $c_n \to A$, akkor
$b_n \to A$.
\[
	c_n
	= \frac{(-32)^n + \frac{1}{(-32)^n}}{16^n + \frac{1}{16^n}}
	= \frac{(-2)^n + \frac{1}{(-512)^n}}{1 + \frac{1}{256^n}}
\]
\[
	c_{2n}
	= \frac{2^{2n} + \frac{1}{512^{2n}}}{1 + \frac{1}{256^{2n}}}
	\to \infty
\]
\[
	c_{2n - 1}
	= \frac{-2^{2n - 1} - \frac{1}{512^{2n - 1}}}{1 + \frac{1}{256^{2n - 1}}}
	\to -\infty
\]
Minden elem szerepel valamelyik részsorozatban, így 2 torlódási pont van:
$\pm \infty$. $\limsup c_n = \infty$, $\liminf c_n = -\infty$, a határérték
nem létezik.
\end{document}
